\chapter{Choosing a published architecture}\label{ch:choice}

\section{The idea of unrolling}

Picture ISTA as a pipeline. One iteration is one box, so $K$ iterations are $K$ boxes in a row. Each box
contains fixed numbers: $D^T/L$ and $\lambda/L$. Now stop treating those numbers as fixed. Make them
trainable, one set per box, and train the whole chain by backpropagation. A loop has become a
neural network.

\begin{center}
\begin{tikzpicture}[node distance=9mm]
\node[fixed] (b) {$b$};
\node[learned,right=of b] (l1) {Layer 1\\ $W_1,\theta_1$};
\node[learned,right=of l1] (l2) {Layer 2\\ $W_2,\theta_2$};
\node[right=of l2] (dots) {$\cdots$};
\node[learned,right=of dots] (lk) {Layer $K$\\ $W_K,\theta_K$};
\node[fixed,right=of lk] (x) {$x$};
\draw[flow] (b)--(l1); \draw[flow] (l1)--(l2); \draw[flow] (l2)--(dots);
\draw[flow] (dots)--(lk); \draw[flow] (lk)--(x);
\end{tikzpicture}
\end{center}

This is \textbf{LISTA}, by Gregor and LeCun (2010). It is written in the ALISTA paper as:
\begin{equation}
x^{(k+1)} = \soft_{\theta^{(k)}}\big(W_1^{(k)} b + W_2^{(k)} x^{(k)}\big). \tag{Liu et al., Eq.~4}
\end{equation}
Setting $W_1=D^T/L$, $W_2=I-D^TD/L$ and $\theta=\lambda/L$ gives back ISTA exactly. So LISTA starts as
ISTA and can only be taught to improve.

\section{Why plain LISTA cannot run here}

$W$ is a full matrix. An image has 1.96 million pixels, so such a matrix would have about
$4\times10^{12}$ entries, and that is for a single pixel-to-pixel map. The earlier dense ALISTA code in
the repository stays as it is, but it cannot run at 0.1\um.

There is a way out. The blur is the same everywhere in the image. So we can use convolutions, which
reuse one small filter at every position. This leads to \textbf{convolutional sparse coding}.

\section{The ALISTA paper}

Liu, Chen, Wang and Yin, ``ALISTA: Analytic Weights Are As Good As Learned Weights in LISTA'', ICLR
2019. It has two main messages.

\begin{enumerate}
\item \textbf{Analysis.} The weights of LISTA do not all need to be learned. A single weight
matrix can be computed from $D$ by a formula. Only step sizes and thresholds then need training.
\item \textbf{Extension.} Section 3 carries all of this to convolutions.
\end{enumerate}

Section 3 starts with the convolutional model, Eq.~17:
\begin{equation}
b = \sum_{m=1}^{M} d_m * x_m^* + \varepsilon . \tag{Liu et al., Eq.~17}
\end{equation}
This is the same equation as our $b\approx\sum_m d_m*x_m$. Then the paper gives the two networks.

\begin{center}
\renewcommand{\arraystretch}{1.5}
\begin{tabularx}{\textwidth}{l>{\raggedright\arraybackslash}X>{\raggedright\arraybackslash}X}
\toprule
& \textbf{Eq.~20: CSC-LISTA} & \textbf{Eq.~26: Conv ALISTA} \\
\midrule
Update &
$x_m \leftarrow \soft_{\theta^{(k)}}\!\big(x_m - (w_m^{(k)})' * (\sum_{\bar m} d_{\bar m}*x_{\bar m} - b)\big)$ &
$x_m \leftarrow \soft_{\theta^{(k)}}\!\big(x_m - \gamma_m^{(k)} \tilde w_m' * (\sum_{\bar m} d_{\bar m}*x_{\bar m} - b)\big)$ \\
Filters & $w_m^{(k)}$ learned, one set per layer & $\tilde w_m$ fixed, computed by formula \\
Trained & All filters and thresholds & Only $\gamma$ and $\theta$ \\
Parameters here & 24.6 million at 32 layers & About a thousand at 32 layers \\
\bottomrule
\end{tabularx}
\end{center}

The prime means a filter rotated by 180 degrees. Think of it as the convolutional version of a
matrix transpose.

\section{Why Eq.~20 and not Eq.~26}

This choice was made during the session, for these reasons:

\begin{enumerate}
\item \textbf{The paper itself reports a small gap.} Its Appendix H denoising test gives Eq.~20 (the paper's ``Conv LISTA'', our CSC-LISTA) an
average PSNR of 31.21 dB and Conv ALISTA 31.07 dB. Learned filters are not worse.
\item \textbf{The analytic filters behaved badly on our blur.} In a small spike, the analytic weight
acted almost like an inverse filter and amplified noise by hundreds of times or more. That spike used
the field itself as the code, with one filter, which is not the model of Chapter~\ref{ch:classics}. Its
code was deleted at your request, so treat this as a note from the session, not a reproducible result.
\item \textbf{You chose.} You told me to drop Eq.~26 and remove everything related to it, so that it
could not cause confusion later. The old dense ALISTA code stayed, also by your choice.
\end{enumerate}

\begin{careful}[The theory does not cover us]
The paper proves fast convergence when the true code is very sparse. For our blur kernel, that limit
worked out to about 1--2 active pixels, while the coil images have tens of thousands or more. So do not claim the
convergence theorems apply here. Use the paper for its architecture, not for its guarantee.
\end{careful}

\section{Other options considered}

\begin{center}
\begin{tabularx}{\textwidth}{lX}
\toprule
Option & Status \\
\midrule
USRNet, ISTA-Net, MoDL & Discussed earlier as published unrolled networks. Not built \\
LFISTA (Moreau and Bruna) & Dense form cannot run. A convolutional relative is Job 3,
Chapter~\ref{ch:lfista} \\
FISTA-Net (Xiang et al.) & You supplied this paper. It is Job 4, Chapter~\ref{ch:fistanet} \\
\bottomrule
\end{tabularx}
\end{center}

\begin{defend}
\textbf{``Where does this network come from?''} The layer is Eq.~20 of Liu et al., ICLR 2019, and
convolutional sparse coding is an established field. What is ours is using it for deblurring, by
writing $d_m=h*g_m$. The paper's own experiment is denoising, so no published result exists for this exact
use. Say this before you are asked.
\end{defend}
